\section{Problema de localización y secuencia metodológica}
\label{bre:fund:problema}

Sea $F:\Omega\subset\R^n\to\R^n$ un campo definido en un abierto.
Se consideran el problema entero $\dot X=F(X)$ y el problema conmensurable
\begin{equation}
 \Caputo q X(t)=F(X(t)),\qquad X(t_0)=X_0,
 \quad 0<q<1.
 \label{bre:fund:pvi}
\end{equation}
La formulación conmensurable asigna el mismo orden a todas las componentes.
El caso $q=1$ se interpreta como una ecuación diferencial ordinaria (EDO).
La estructura tipo Chua y el orden de la derivada son características
independientes: una realización de Lur'e puede describir un sistema no Chua.

En orden entero, sea $\varphi_t$ el flujo hacia adelante, sobre los datos
cuyas soluciones existen. Para un compacto invariante y atractivo $A$, su cuenca es
\begin{equation}
 \mathcal B(A)=\{X_0:\operatorname{dist}(\varphi_t(X_0),A)\to0\}.
 \label{bre:fund:cuenca}
\end{equation}
Aquí se exige existencia global de las trayectorias incluidas en la cuenca.
La clasificación usual denomina oculto a un atractor cuya cuenca no
interseca vecindades suficientemente pequeñas de ningún equilibrio
\cite{LeonovKuznetsov2013,DudkowskiEtAl2016}. Con
$\mathcal E=\{e:F(e)=0\}$, la condición precisa es
\begin{equation}
 \forall e\in\mathcal E\ \exists r_e>0:
 B(e,r_e)\cap\mathcal B(A)=\varnothing.
 \label{bre:fund:hidden}
\end{equation}
El orden de los cuantificadores permite radios distintos para equilibrios
distintos. Un contacto a gran distancia de $e$ no contradice la existencia
de un radio menor excluido de la cuenca. Si $\mathcal E$ es vacío, la
condición es vacua; todavía debe establecerse la existencia de un atractor.

La metodología comprende seis operaciones encadenadas:
\begin{enumerate}
 \item Fijar campo, parámetros, dominio, orden, terminal inferior y
 convención de solución; resolver los equilibrios y su estabilidad local.
 \item Construir semillas mediante balance armónico, puntos perpetuos,
 continuación desde un problema resoluble o búsqueda acotada.
 \item Transportar las semillas al campo objetivo conservando la historia
 en los problemas fraccionarios.
 \item Integrar el campo autónomo bajo un mismo convenio de inicialización,
 separar el transitorio y comparar paso, horizonte y memoria.
 \item Identificar destinos mediante trayectorias, secciones, espectros,
 recurrencia e indicadores variacionales compatibles con el modelo.
 \item Estudiar la atracción y la separación respecto de los equilibrios;
 aplicar retornos, regiones invariantes o métodos topológicos cuando sus
 hipótesis estén verificadas.
\end{enumerate}
La resolución algebraica determina candidatos; la integración determina
su destino a tiempo finito. La demostración de atracción y la exclusión de
vecindades completas son problemas adicionales, con hipótesis propias.

\section{Fundamento analítico de la formulación fraccionaria}
\label{bre:fund:caputo}

\subsection{Integral fraccionaria y dominio del operador}

Para $a>0$ y $g\in L^\infty([t_0,T],\R^n)$, se define
\begin{equation}
 (I_{t_0}^a g)(t)=\frac1{\Gamma(a)}
 \int_{t_0}^{t}(t-s)^{a-1}g(s)\dd s.
 \label{bre:fund:integral}
\end{equation}
El núcleo es integrable en $s=t$ porque $a-1>-1$. En particular,
\begin{equation}
 \norm{I_{t_0}^a g(t)}\le
 \frac{(t-t_0)^a}{\Gamma(a+1)}\norm g_\infty,
 \label{bre:fund:integralbound}
\end{equation}
al integrar la potencia y usar $\Gamma(a+1)=a\Gamma(a)$.
Esta cota prueba la continuidad en el terminal inferior y la anulación
del incremento allí.

La composición $I^aI^b=I^{a+b}$ requiere el mismo terminal inferior.
Para justificarla, escriba la integral doble sobre
$t_0\le s\le u\le t$. Su integral absoluta es finita por la acotación de
$g$ y la integrabilidad de ambos núcleos; Fubini permite cambiar el orden.
La integral interior se calcula con $u=s+(t-s)v$:
\begin{align}
 &\int_s^t(t-u)^{a-1}(u-s)^{b-1}\dd u \nonumber\\
 &\quad=(t-s)^{a+b-1}\int_0^1(1-v)^{a-1}v^{b-1}\dd v\nonumber\\
 &\quad=(t-s)^{a+b-1}\frac{\Gamma(a)\Gamma(b)}{\Gamma(a+b)}.
 \label{bre:fund:beta}
\end{align}
Sustituir esta expresión en la integral exterior da la composición.
La identidad se extiende a $g\in L^1$: Tonelli aplicado al valor absoluto
proporciona, con $H=T-t_0$,
\[
 \norm{I^a g}_{L^1}\le\frac{H^a}{\Gamma(a+1)}\norm g_{L^1}.
\]
La integral doble es entonces absolutamente integrable en el intervalo y
Fubini demuestra la composición casi en todas partes. Esta extensión
permite aplicarla a la derivada ordinaria de una función absolutamente continua.
La misma sustitución demuestra, para $\mu>-1$,
\begin{equation}
 I_{t_0}^a(t-t_0)^\mu
 =\frac{\Gamma(\mu+1)}{\Gamma(\mu+a+1)}(t-t_0)^{\mu+a}.
 \label{bre:fund:power}
\end{equation}

Si $X$ es absolutamente continua, la derivada clásica de Caputo es
\begin{equation}
 \Caputo q X=I_{t_0}^{1-q}X'
 =\frac1{\Gamma(1-q)}\int_{t_0}^{t}(t-s)^{-q}X'(s)\dd s,
 \label{bre:fund:classical}
\end{equation}
donde la integral se entiende en los puntos en los cuales existe, al menos
casi en todas partes. Para no imponer una regularidad temporal que no se
haya demostrado, se utiliza también la definición
\begin{equation}
 \Caputo q X=\frac{\dd}{\dd t}I_{t_0}^{1-q}(X-X_0),
 \label{bre:fund:generalized}
\end{equation}
cuando $X$ es continua, $X(t_0)=X_0$ y el término derivado es
absolutamente continuo. Ambas definiciones coinciden sobre su dominio
común: $X-X_0=I^1X'$ y la composición de integrales da
$I^{1-q}(X-X_0)=I^1I^{1-q}X'$, cuya derivada es $I^{1-q}X'$.
Las soluciones se buscan en la clase $X_0+I_{t_0}^q(L^\infty)$ sobre
cada intervalo compacto de existencia \cite{Gomoyunov2018Convex}.

\subsection{Elección de Caputo y significado del dato inicial}
\label{bre:fund:initialchoice}

La elección del operador forma parte del modelo. Para $0<q<1$, Caputo
permite prescribir el valor ordinario $X(t_0)=X_0$ y anula las constantes.
La derivada de Riemann--Liouville, definida por
$D_{\rm RL}^qX=(\dd/\dd t)I^{1-q}X$, satisface sobre el dominio anterior
\begin{equation}
 D_{\rm RL}^qX=\Caputo qX+
 \frac{X_0}{\Gamma(1-q)}(t-t_0)^{-q},\quad t>t_0.
 \label{bre:fund:RLbridge}
\end{equation}
En efecto, se separa $X=(X-X_0)+X_0$ y se deriva
$I^{1-q}X_0=X_0(t-t_0)^{1-q}/\Gamma(2-q)$.
Por tanto, el mismo PVI puede escribirse con Riemann--Liouville si se
incorpora ese término singular de inicialización. Sustituir el operador
y conservar sin cambios el segundo miembro produce, en general, otro
problema. La conveniencia de Caputo aquí procede de los datos iniciales
puntuales y de la formulación causal elegida; no establece una superioridad
universal sobre otros operadores.

El valor inicial de $X$ no debe confundirse con el valor de su derivada
fraccionaria en el extremo. Para el campo constante $F(X)=v$, Volterra da
\begin{equation}
 X(t)=X_0+\frac{(t-t_0)^q}{\Gamma(q+1)}v.
 \label{bre:fund:startupconstant}
\end{equation}
Su derivada ordinaria es $v(t-t_0)^{q-1}/\Gamma(q)$, integrable pero
no acotada cuando $v\ne0$. Sustituirla en Caputo y aplicar la integral beta
da $\Caputo qX(t)=v$ para todo $t>t_0$. Así, el límite por la derecha
puede ser distinto de cero. Asignar cero a una integral de intervalo vacío
en $t=t_0$ es una convención de extremo; no autoriza a imponer
$F(X_0)=0$. La ecuación se exige para $t>t_0$ en el sentido de solución
declarado, y el dato inicial se obtiene por continuidad de $X$.

El operador de Liouville--Weyl, con historia desde $-\infty$, permite
analizar armónicos estacionarios. Su relación con el operador causal de
Caputo justifica el predictor armónico mediante un término explícito de
historia; no identifica ambos problemas iniciales. Las diferencias de
Grünwald--Letnikov pueden aproximar operadores compatibles con estas
formulaciones cuando se fijan la inicialización y la memoria apropiadas.
El nombre de la discretización, por sí solo, no determina el PVI que resuelve.

\subsection{Equivalencia con Volterra y existencia local}

Si $g(t)=F(X(t))$ es acotada, la ecuación integral correspondiente es
\begin{equation}
 X(t)=X_0+I_{t_0}^q[F(X)](t).
 \label{bre:fund:volterra}
\end{equation}
Aplicar $I^{1-q}$ a $X-X_0$ y usar \eqref{bre:fund:beta} produce
$I^1g$. Su derivada es $g$ casi en todas partes, de modo que
\eqref{bre:fund:generalized} satisface \eqref{bre:fund:pvi}.
Recíprocamente, integrar \eqref{bre:fund:generalized} da
$I^{1-q}(X-X_0)=I^1g$: la constante es cero por la continuidad en $t_0$.
Al restar $I^{1-q}I^qg=I^1g$, queda $I^{1-q}h=0$, con
$h=X-X_0-I^qg$. Aplicar $I^q$ da $I^1h=0$; al derivar se obtiene
$h=0$ casi en todas partes y, por continuidad, en todo el intervalo.
Así se demuestra la equivalencia, sin introducir $X'$ donde no esté
justificado \cite{Podlubny1999}.

Supóngase ahora $F$ localmente Lipschitz en $X$. Elija una bola cerrada
$\overline B(X_0,R)\subset\Omega$ y constantes $M,L$ tales que
$\norm{F(X)}\le M$ y
$\norm{F(X)-F(Y)}\le L\norm{X-Y}$ en ella. Sobre
\[
 \mathcal X=\{X\in C([t_0,t_0+\tau],\R^n):
               \norm{X-X_0}_\infty\le R\}
\]
se define $\mathcal TX=X_0+I^q[F(X)]$. La estimación anterior da
\begin{align}
 \norm{\mathcal TX-X_0}_\infty&\le\frac{M\tau^q}{\Gamma(q+1)},\nonumber\\
 \norm{\mathcal TX-\mathcal TY}_\infty
 &\le\frac{L\tau^q}{\Gamma(q+1)}\norm{X-Y}_\infty.
 \label{bre:fund:contraction}
\end{align}
Por tanto, escoger $M\tau^q/\Gamma(q+1)\le R$ y
$L\tau^q/\Gamma(q+1)<1$ hace de $\mathcal T$ una contracción de un
subconjunto cerrado de un espacio de Banach en sí mismo. Su único punto
fijo es la solución local. Para campos dependientes del tiempo se requieren
las mismas cotas uniformes, además de medibilidad temporal y continuidad
espacial; la formulación integral cubre campos por tramos durante la
continuación. Un campo PWL continuo es localmente Lipschitz, aunque su
jacobiano no exista en las superficies de conmutación.

La extensión a tiempos mayores conserva la integral acumulada. Si la
solución permanece en un compacto de $\Omega$ hasta un supuesto tiempo
máximo finito, $F(X)$ permanece acotada. La cota de continuidad de $I^q$
proporciona un límite al aproximarse a ese tiempo y la misma construcción
contractiva, con la historia fijada, extiende la solución. En consecuencia,
la pérdida de existencia global exige abandonar todos los compactos del
dominio. Esta alternativa se usa al demostrar acotación en Lorenz.

\subsection{Transformada de Laplace y respuesta lineal}

Trasladando $t_0$ a cero, supóngase $X-X_0=I^qg$ y
$e^{-\sigma t}g\in L^1([0,\infty))$ para algún $\sigma>0$.
Fubini con peso exponencial garantiza las transformadas de $X-X_0$,
$H=I^{1-q}(X-X_0)=I^1g$ y $H'=g$ en un semiplano derecho.
La transformada del núcleo de \eqref{bre:fund:integral}
es $s^{-a}$. Se obtiene de la integral gamma mediante $v=st$ para
$s>0$ real; la identidad se extiende analíticamente al semiplano de
convergencia con la rama principal. Fubini, ahora con peso exponencial,
convierte la convolución en producto. Como
$I^{1-q}(X-X_0)$ se anula inicialmente,
\begin{align}
 \mathcal L\{\Caputo qX\}(s)
 &=s\,s^{q-1}\left(\widehat X(s)-\frac{X_0}{s}\right)\nonumber\\
 &=s^q\widehat X(s)-s^{q-1}X_0.
 \label{bre:fund:laplace}
\end{align}
El término inicial no desaparece salvo cuando $X_0=0$.
Para $\Caputo qX=AX+Bu$, $y=CX$, se sigue
\begin{align}
 \widehat X(s)&=(s^qI-A)^{-1}s^{q-1}X_0\nonumber\\
 &\quad+(s^qI-A)^{-1}B\widehat u(s),\nonumber\\
 W_q(s)&=C(s^qI-A)^{-1}B.
 \label{bre:fund:transfer}
\end{align}
La transferencia describe la respuesta a la entrada con condición inicial
nula. En el balance armónico se usa
$(\ii k\omega)^q=(k\omega)^q e^{\ii q\pi/2}$ para $k>0$ y el conjugado
para $k<0$. El modo constante tiene multiplicador cero. Se fija
$s^q=\exp(q\operatorname{Log}s)$, con
$\operatorname{Arg}s\in(-\pi,\pi)$.

La expansión del resolvente para $\norm A<|s|^q$ produce
\[
 (s^qI-A)^{-1}s^{q-1}
 =\sum_{j=0}^\infty A^j s^{-qj-1}.
\]
La transformada inversa de cada término y la convergencia de la serie
de Mittag--Leffler dan
\begin{equation}
 X(t)=E_q(At^q)X_0,\quad
 E_{q,\beta}(Z)=\sum_{j=0}^\infty\frac{Z^j}{\Gamma(qj+\beta)},
 \label{bre:fund:ML}
\end{equation}
con $E_q=E_{q,1}$. Para una fuerza $g$, el mismo resolvente da
\begin{equation}
 X(t)=E_q(At^q)X_0+
 \int_0^t K_A(t-s)g(s)\dd s,
 \label{bre:fund:variation}
\end{equation}
donde $K_A(t)=t^{q-1}E_{q,q}(At^q)$. La convolución se interpreta sobre
intervalos finitos aun cuando no se haya demostrado una cota exponencial
global para el problema no lineal.
Para justificar la fórmula sin una inversión formal, supóngase
$g\in L^\infty([0,T])$ e itérese
$X=X_0+I^q(AX+g)$. La composición de integrales produce las series
$\sum_{j\ge0}A^jt^{qj}X_0/\Gamma(qj+1)$ y
$\sum_{j\ge0}A^jI^{q(j+1)}g$. Sus normas están dominadas,
respectivamente, por las series escalares con términos
$\norm A^jT^{qj}\norm{X_0}/\Gamma(qj+1)$ y
$\norm A^jT^{q(j+1)}\norm g_\infty/\Gamma(q(j+1)+1)$.
Ambas convergen; la convergencia uniforme permite integrar la suma,
verificar Volterra y obtener \eqref{bre:fund:variation} por unicidad.

\subsection{Equilibrios, linealización y estabilidad}

Un equilibrio satisface $F(e)=0$, porque Caputo anula las constantes.
Si $F\in C^1$ cerca de $e$, escriba
\begin{equation}
 F(e+\xi)=J_e\xi+R(\xi),\qquad J_e=DF(e).
 \label{bre:fund:linearization}
\end{equation}
El teorema fundamental del cálculo en una bola convexa proporciona
\begin{align}
 R(\xi)-R(\eta)
 &=\int_0^1[DF(e+\eta+s(\xi-\eta))-J_e]\nonumber\\
 &\hspace{24mm}\cdot(\xi-\eta)\dd s,\nonumber\\
 \norm{R(\xi)-R(\eta)}&\le\ell_r\norm{\xi-\eta},\nonumber\\
 \ell_r&=\sup_{\norm u\le r}\norm{DF(e+u)-J_e}\longrightarrow0.
 \label{bre:fund:smallremainder}
\end{align}
En particular, $R(\xi)=o(\norm\xi)$; una cota cuadrática requiere
$DF$ localmente Lipschitz, por ejemplo $F\in C^2$. En un campo PWL,
este argumento clásico se aplica a equilibrios interiores a una región.

El criterio de Matignon para $0<q<1$ establece estabilidad asintótica del
sistema lineal si y sólo si todos los autovalores son no nulos y
\begin{equation}
 |\operatorname{Arg}\lambda_j(J_e)|>q\pi/2.
 \label{bre:fund:matignon}
\end{equation}
La condición caracteriza el decaimiento de \eqref{bre:fund:ML}, incluidos
los bloques de Jordan \cite{Matignon1996}. Su aplicación no lineal requiere
el teorema de linealización y la pequeñez Lipschitz del resto
\cite{CongDoanSiegmundTuan2016}. El mecanismo puede verse en
\eqref{bre:fund:variation}: bajo el sector estricto existen
\begin{align*}
 M_A&=\sup_{t\ge0}\norm{E_q(J_et^q)}<\infty,\\
 C_A&=\int_0^\infty\norm{K_{J_e}(s)}\dd s<\infty.
\end{align*}
Estas estimaciones son consecuencias de las cotas sectoriales de
Mittag--Leffler. Cerca de cero el núcleo es del orden $t^{q-1}$,
integrable; en el infinito su decaimiento es del orden $t^{-q-1}$.
Se elige $r$ con $C_A\ell_r<1$ y
$M_A\norm{\xi_0}<(1-C_A\ell_r)r$.
La fórmula de variación impide que la solución alcance por primera vez
la frontera de la bola y da estabilidad. Para la atracción, se separa
la convolución en una parte de tiempo inicial fijo, que tiende a cero,
y una cola. Si $L=\limsup_{t\to\infty}\norm{\xi(t)}$, resulta
$L\le C_A\ell_rL$, de donde $L=0$.
Así se explicita la función de cada hipótesis del resultado no lineal.

El margen $m_q(e)=\min_j(|\operatorname{Arg}\lambda_j|-q\pi/2)$
cuantifica la distancia angular a la frontera. $m_q>0$ verifica el criterio
estricto; $m_q<0$ señala un modo lineal creciente. Si $m_q=0$ o existe un
autovalor nulo, el argumento anterior no decide la estabilidad. Para $q=1$
se recupera el criterio de Hurwitz $\operatorname{Re}\lambda_j<0$.
Un margen pequeño requiere controlar el error de los autovalores antes
de asignar su signo.

\subsection{Memoria y espacio de estados}

La evolución puntual de Caputo no satisface, en general, la ley de
semigrupo. El ejemplo escalar $\Caputo q x=\lambda x$ lo muestra:
\begin{align}
 &E_q\bigl(\lambda(2h)^q\bigr)-E_q(\lambda h^q)^2\nonumber\\
 &\quad=\frac{\lambda(2^q-2)}{\Gamma(1+q)}h^q+O(h^{2q}).
 \label{bre:fund:nosemigroup}
\end{align}
Para $\lambda\ne0$ y $0<q<1$ el coeficiente principal no es cero.
La composición de integrales de \eqref{bre:fund:beta} concierne a los
órdenes de los operadores; no autoriza a reiniciar la evolución temporal.

Si $t_0=0$, dividir Volterra en $[0,t]$ y $[t,t+\theta]$ produce
\begin{align}
 m_t(\theta)&=X_0+\frac1{\Gamma(q)}
  \int_0^t(t+\theta-s)^{q-1}F(X(s))\dd s,\nonumber\\
 X(t+\theta)&=m_t(\theta)+\frac1{\Gamma(q)}
  \int_0^\theta(\theta-u)^{q-1}F(X(t+u))\dd u.
 \label{bre:fund:profile}
\end{align}
El perfil $m_t$ determina la continuación; $m_t(0)=X(t)$ es sólo una
evaluación. La construcción de un semiflujo sobre perfiles exige
bien planteamiento global y continuidad en el espacio elegido
\cite{DoanKloeden2021Semidynamical}. No basta adoptar sin restricciones
todo $C([0,\infty),\R^n)$ con la topología de convergencia uniforme en
compactos. En efecto, para cualquier $c$, el término libre
\begin{equation}
 m_c(\theta)=c-\frac{\theta^q}{\Gamma(q+1)}F(c)
 \label{bre:fund:artificial}
\end{equation}
produce la solución constante $X\equiv c$. La continuación conserva
$m_c$, aunque $F(c)\ne0$. En el campo $F(x)=-x$, estos perfiles fijos
se aproximan al perfil cero cuando $c\to0$: el equilibrio estable de
reinicio no atrae una vecindad de todos los perfiles.

Para definir atracción funcional se especifica un espacio admisible
$\mathcal X_{\rm adm}$, positivamente invariante, con semiflujo continuo
$S_t$. Para un compacto estrictamente invariante $A$ y una familia
de reinicios $\iota(B)$, con $\iota(X_0)(\theta)=X_0$, la atracción uniforme
requiere $\iota(B)\subset\mathcal X_{\rm adm}$ y
significa
\begin{equation}
 \sup_{X_0\in B}\operatorname{dist}_{\mathcal X}
       (S_t\iota(X_0),A)\longrightarrow0.
 \label{bre:fund:resetattraction}
\end{equation}
Se debe indicar qué conjuntos $B$ se incluyen; un atractor local no tiene
que atraer todos los reinicios acotados. La ocultedad funcional excluye
vecindades de los perfiles de equilibrio en $\mathcal X_{\rm adm}$.
La ocultedad bajo reinicio excluye bolas de $X_0$ según
\eqref{bre:fund:hidden}, usando esa misma atracción. Los sondeos
numéricos de este estudio exploran reinicios finitos y destinos observables;
no construyen por sí solos un espacio funcional admisible ni verifican
la convergencia uniforme de \eqref{bre:fund:resetattraction}.

\subsection{Acotación sin aplicar una regla de la cadena ordinaria}

Para Caputo no se sustituye $\Caputo q[V(X)]$ por
$\nabla V(X)\cdot F(X)$. Si $V$ es convexa y $C^1$ con gradiente
localmente Lipschitz, y $X-X_0\in I^q(L^\infty)$, se dispone de
\begin{equation}
 \Caputo q[V(X(t))]\le\nabla V(X(t))\cdot\Caputo qX(t).
 \label{bre:fund:convex}
\end{equation}
La desigualdad se cumple casi en todas partes. El resultado también asegura la regularidad fraccionaria de
$V(X)-V(X_0)$ \cite{Gomoyunov2018Convex}. Para ver el signo en el caso
suave, fije $t$ y defina el resto convexo
$\Phi(s)=V(X(s))-V(X(t))-\nabla V(X(t))\cdot(X(s)-X(t))\ge0$.
Entonces $\Phi(t)=0$ y la diferencia entre ambos miembros, multiplicada
por $\Gamma(1-q)$, es
\begin{align}
 \int_0^t(t-s)^{-q}\Phi'(s)\dd s
 &=-t^{-q}\Phi(0)\nonumber\\
 &\quad-q\int_0^t(t-s)^{-q-1}\Phi(s)\dd s\le0.
 \label{bre:fund:convexproof}
\end{align}
La integración por partes se hace primero hasta $t-\varepsilon$.
Con trayectoria Lipschitz y gradiente Lipschitz,
$\Phi(s)=O((t-s)^2)$, por lo cual el término del extremo tiende a cero
y la integral converge. Para la clase $I^q(L^\infty)$, la extensión se
realiza por aproximación y compacidad; no se presupone la existencia de
$X'$ en toda esa clase. En la sección de resultados se aplica esta
desigualdad al Lorenz generalizado y se exhibe la cota obtenida.

Los modelos autónomos de Caputo con terminal inferior fijo y orden no
entero no admiten las órbitas periódicas no constantes propias de una EDO
bajo las hipótesis del teorema de no periodicidad
\cite{TavazoeiHaeri2009,AreaLosadaNieto2014}. Por ello, un ciclo del
problema armónico estacionario, una curva cerrada dibujada por una muestra
y una solución periódica exacta del PVI causal son objetos distintos.
La construcción armónica emplea el primero como predictor; la evaluación
numérica posterior se refiere siempre al PVI declarado.

\subsection{Qué permite trasladar un método de orden entero}
\label{bre:fund:transfermethods}

Las ecuaciones $F(e)=0$ y las realizaciones algebraicas del mismo campo no
dependen de $q$. Su reutilización es exacta. En cambio, sustituir
$\ii\omega$ por $(\ii\omega)^q$ sólo traslada la respuesta del operador
lineal estacionario: la selección de una trayectoria causal requiere
resolver nuevamente Volterra. La misma separación se aplica a los puntos
perpetuos, cuya ecuación algebraica se conserva, pero cambia su
interpretación como aceleración.

Una relación rigurosa con el caso entero puede establecerse a tiempo
finito. Sean $X_q$ y $X_1$ soluciones con el mismo dato inicial, sobre
$[0,T]$, contenidas en un compacto común en el cual $F$ esté acotado por
$M$ y sea Lipschitz con constante $L$. Ponga
\[
 k_q(t)=\frac{t^{q-1}}{\Gamma(q)},\qquad
 \varepsilon_q=\int_0^T|k_q(t)-1|\dd t.
\]
Para $q\in[q_*,1]$, $q_*>0$, el integrando está dominado por un múltiplo
de $t^{q_*-1}+1$, integrable en $[0,T]$. Como $k_q(t)\to1$ para
$t>0$ cuando $q\to1$, la convergencia dominada da
$\varepsilon_q\to0$. Con $e_q(t)=\norm{X_q(t)-X_1(t)}$, restar ambas
ecuaciones integrales produce
\begin{align}
 e_q(t)
 &\le M\varepsilon_q\nonumber\\
 &\quad+L\int_0^t k_q(t-s)e_q(s)\dd s,\nonumber\\
 \sup_{0\le t\le T}e_q(t)
 &\le M\varepsilon_q E_q(LT^q)\longrightarrow0.
 \label{bre:fund:qtoone}
\end{align}
La última desigualdad procede de iterar la integral positiva:
el término de orden $j$ es
$(Lt^q)^j/\Gamma(qj+1)$, y su suma es $E_q(Lt^q)$.
El factor de Mittag--Leffler permanece uniformemente acotado para
$q\in[q_*,1]$: ponga $B=L\max(T,T^{q_*})$. Para $j$ suficientemente
grande, la monotonía de $\Gamma$ en el semieje correspondiente permite
dominar su término por $B^j/\Gamma(q_*j+1)$; la serie dominante converge.
Los términos restantes son un número finito de funciones continuas de $q$.
La hipótesis de compacto común se verifica directamente o se utiliza la
estimación hasta el primer tiempo de salida. Esta prueba justifica
comparar trayectorias enteras y fraccionarias en una ventana controlada.
No proporciona una estimación uniforme para $T\to\infty$; por tanto, no
demuestra persistencia de un atractor, de su cuenca ni de su ocultedad al
variar $q$. Esos problemas requieren controlar por separado la dinámica
asintótica y el espacio de estados.

\begin{table*}[t]
\centering
\caption{Condiciones de adaptación de los métodos enteros al problema causal de Caputo.}
\label{bre:fund:adaptationtable}
\small
\begin{tabularx}{\textwidth}{@{}L{0.19\textwidth}L{0.30\textwidth}Y@{}}
\toprule
Operación & Resultado utilizable & Hipótesis o desarrollo adicional\\
\midrule
Equilibrios y realización & $F(e)=0$ e identidad de la realización del campo. &
Mismo campo, dominio y parámetros; tratar las regiones no suaves antes de derivar.\\
Linealización & Matriz $DF(e)$ y resto de Taylor. &
Regularidad $C^1$, pequeñez Lipschitz del resto y sector estricto de Matignon.\\
Balance armónico & Multiplicador de Fourier $(\ii k\omega)^q$ y coeficientes de la no linealidad. &
Problema auxiliar estacionario; rama compleja fijada, residuo físico y control de armónicos omitidos.\\
Continuación & Familia de campos con extremos algebraicamente comprobados. &
Una sola historia para la familia causal; un reinicio posterior define un nuevo PVI. La persistencia asintótica no se sigue de continuidad a tiempo finito.\\
Puntos perpetuos & Ecuación $DF\,F=0$ para generar semillas. &
Expansión de arranque en $t^q$; no identificarla con $X''=0$. Los eventos secuenciales dependen de la historia.\\
Retornos y cuencas & Cruces observados y destinos de reinicios. &
Operador funcional cuando se conserva memoria; hipótesis de retorno, atracción y exclusión de vecindades completas.\\
Lyapunov y compacidad & Desigualdad convexa y estimaciones integrales. &
Verificar la clase $I^q(L^\infty)$, continuación global y el espacio de perfiles; no usar una regla de la cadena ordinaria.\\
\bottomrule
\end{tabularx}
\end{table*}

Danca estudia Lorenz generalizado, Rabinovich--Fabrikant y Chua PWL con
Caputo, estabilidad sectorial y un predictor--corrector fraccionario
\cite{Danca2017}. Wu et al. construyen condiciones iniciales para Chua
con arctangente mediante una adaptación del procedimiento de
Leonov y emplean ADM en la simulación \cite{Wu2023ArctanChua}.
Ambos trabajos reconocen la dificultad de interpretar la periodicidad en
orden fraccionario; Wu et al. también discuten la pérdida de historia de
su aproximación por intervalos. Estas contribuciones suministran modelos,
parámetros y procedimientos de búsqueda. La validez de una operación
concreta se determina, además, por las hipótesis de la
Tabla~\ref{bre:fund:adaptationtable}. Reproducir una recurrencia publicada
y aproximar el PVI de Caputo con memoria completa son verificaciones
distintas, que se comparan mediante el operador integral correspondiente.
